\vspace{2.5pt plus 3pt minus 0.5pt}\noindent\hspace*{1.2em}\kbvnum{(8)}\hspace{0.55em}ತದ್ದಿತಾಂತಾವ್ಯಯಗಳೆಂದರೇನು? ಉದಾಹರಣೆ ಸಹಿತ ವಿವರಿಸಿ.\par

\vspace{2.5pt plus 3pt minus 0.5pt}\noindent\hspace*{1.2em}\kbvnum{(9)}\hspace{0.55em}ಈ ಅರ್ಥ ಬರುವ ಒಂಟಿಪದ ಯಾವುದು ?\par

\vspace{2.5pt plus 3pt minus 0.5pt}\noindent\hspace*{1.2em}\kbvnum{(ಅ)}\hspace{0.55em}ಬಳೆಗಳನ್ನು ಮಾರುವವನು (ಬ) ಸುಂದರ ರೂಪ ಉಳ್ಳವಳು\par

\vspace{2.5pt plus 3pt minus 0.5pt}\noindent\hspace*{1.2em}\kbvnum{(ಕ)}\hspace{0.55em}ಸಿರಿ ಉಳ್ಳವನ ಭಾವ (ಡ) ಗೊಲ್ಲರ ಹೆಂಗಸು\par

\vspace{2.5pt plus 3pt minus 0.5pt}\noindent\hspace*{1.2em}\kbvnum{(10)}\hspace{0.55em}ಈ ಪದಗಳಲ್ಲಿರುವ ತದ್ದಿತ ಪ್ರತ್ಯಯಗಳನ್ನು ತಿಳಿಸಿ : ಹಾವಾಡಿಗ, ಒಕ್ಕಲಗಿತ್ತಿ, ಜಾಣತನ, ರಾಮನಂತೆ ತಾಯಿಗಿಂತ\par

\kbvchapter{ಪಾಠ ೧೬}{ಅವ್ಯಯಗಳು}

\kbvsechead{ಈ ವಾಕ್ಯಗಳನ್ನು ಓದಿ ನೋಡೋಣ :}

\noindent\hspace*{2.2em}ಗಿಡವು ಚೆನ್ನಾಗಿ ಬೆಳೆಯುತ್ತಿದೆ.\par

\indent ಜ್ಯೋತಿಯು ಚೆನ್ನಾಗಿ ಹಾಡುತ್ತಾಳೆ. ಅಕ್ಕಾ ! ಚೆನ್ನಾಗಿ ಆಲೋಚಿಸು, ನದಿಯಲ್ಲಿ ಚೆನ್ನಾಗಿ ಈಜಬಹುದು.\par

\noindent\hspace*{2.2em}ಹುಡುಗರು ಚೆನ್ನಾಗಿ ಓಡುತ್ತಿದ್ದಾರೆ.\par

\noindent\hspace*{2.2em}ಅವರಿಗೆ ಚೆನ್ನಾಗಿ ಹಣ ಸಿಕ್ಕಿತು.\par

\noindent\hspace*{2.2em}ಮಕ್ಕಳು ಚೆನ್ನಾಗಿ ಆಟವಾಡಿದರು.\par

\noindent\hspace*{2.2em}ಪಕ್ಷಿಗಳು ಚೆನ್ನಾಗಿ ಹಾರುತ್ತಿವೆ.\par

\noindent\hspace*{2.2em}ಮರಗಳಲ್ಲಿ ಚೆನ್ನಾಗಿ ಹೂಬಿಟ್ಟಿದೆ.\par

\indent ಮೇಲಿನ ವಾಕ್ಯಗಳಲ್ಲಿ ಪ್ರತಿಯೊಂದರಲ್ಲೂ ಚೆನ್ನಾಗಿ ಎಂಬ ಶಬ್ದ ಬಂದಿದೆ ; ನಾಮಪದ ಕ್ರಿಯಾಪದಗಳಲ್ಲಿ ಬದಲಾವಣೆಯಾಗಿದ್ದರೂ “ಚೆನ್ನಾಗಿ” ಎಂಬ ಶಬ್ದ ಇದ್ದಂತೆಯೇ ಇರುವುದು. ಹೀಗೆ ಲಿಂಗ ವಚನ ವಿಭಕ್ತಿಗಳಲ್ಲಿ ಬದಲಾವಣೆಯುಂಟಾದರೂ ಸ್ವತಃ ಯಾವುದೇ ವ್ಯತ್ಯಾಸ ಹೊಂದದ ಶಬ್ದವು ಅವ್ಯಯವೆನಿಸುವುದು.\par

\indent ಸುಲಭವಾಗಿ, ತಟ್ಟನೆ, ಭರ್ರನೆ, ಅಯ್ಯೋ !, ಬೇಕು, ಮತ್ತು ಅಥವಾ, ಅಂತೆ, ಬರಲು, ಛೀ !, ಧಗಧಗ—ಇವು ಅವ್ಯಯಗಳಿಗೆ ಕೆಲವು ಉದಾಹರಣೆಗಳು.\par
